\section{抓包分析：API 层面的理解}

讲者通过 HTTP 请求/响应的 Hook 对 Gemini API 进行了抓包分析：

\begin{itemize}
    \item \textbf{Request Body}：包含 Tools 定义（Code Execution）和用户输入（图片 + 文本）
    \item \textbf{Response Body}：包含多个 Parts，每个 Part 可能是：
    \begin{itemize}
        \item \texttt{executable\_code}：模型生成的 Python 代码
        \item \texttt{code\_execution\_result}：代码执行的结果（包括生成的图片）
        \item \texttt{text}：模型的文本推理输出
    \end{itemize}
    \item 前端展示的每一步过程，本质上就是在可视化这些 Parts
\end{itemize}

\begin{knowledgebox}{Gemini App vs AI Studio}
截至视频录制时，Gemini App（消费者版本）尚未提供 Code Execution Tool。要体验 Agentic Vision，需要使用 AI Studio（开发者工具）。
\end{knowledgebox}


